\documentclass[shownotes]{beamer}
\input{../../aux/diapo_encabezado.tex}

\title[Bayesian inference]{Evidence in favor of a scientific theory\\ \large With great dishonesty comes great overfitting}

\author[Gustavo Landfried]{Gustavo Landfried \\ \scriptsize \texttt{@GALandfried} \\\vspace{0.2cm}
MSc in Anthropological Sciences \\
PhD student in Computer Sciences 
\vspace{-0.3cm}}
\institute[DC-ICC-CONICET]{\includegraphics[width=0.4\textwidth]{../../Imagenes/dc-logo}}
\date{}

\begin{document}
\begin{frame}[noframenumbering]
 
 \begin{textblock}{100}(23,18)
 \includegraphics[width=0.3\textwidth]{../../Imagenes/logo_licar} 
 \end{textblock}
  \begin{textblock}{100}(74,18)
 \includegraphics[width=0.3\textwidth]{../../Imagenes/LIAA/logo_version_02} 
 \end{textblock}

\vspace{2.5cm}
\maketitle
 
\end{frame}

\footnotesize

\section{Conditional probability}


\begin{frame}
\begin{textblock}{128}(0,8)
\begin{center}
 \Large Conditional probability
\end{center}
\end{textblock}
\vspace{0.75cm}


\begin{table}[H]
\begin{tabular}{c|c|c|c}
 & $A_1$ & $A_2$ &  \\ \hline
 $B_1$& 4 & 2 & 6 \\ \hline
 $B_2$& 76 & 18 & 94\\ \hline
 & 80 & 20 & 100
\end{tabular}
\end{table}

\begin{center}
What is the chance of an $A_1$ event given $B_1$? 
\end{center}

\pause

\begin{equation*}
 P(A_1|B_1) = \frac{P(B_1 \cap A_1)}{P(B_1)}
\end{equation*}

\pause 

\begin{mdframed}
Bayes theorem:
\begin{equation}
  P(A_1|B_1) = \frac{P(B_1|A_1)P(A_1)}{P(B_1)}
\end{equation}
\end{mdframed}

\end{frame}

\subsection{Scientific test example}

\begin{frame}
\begin{textblock}{128}(0,8)
\begin{center}
 \large Scientific test example
\end{center}
\end{textblock}
\vspace{0.5cm}


{ \footnotesize
 There is a test that correctly detects vampirism $95\%$ of the time.
 
 $\bullet$ $P(\text{positive}|\text{vampire})=0.95$
 
 \pause \medskip 

 One percent of the time it incorrectly detect normal personas as vampires.
 
 $\bullet$ $P(\text{positive}|\text{mortal})=0.01$
 
 \pause \medskip 
 
 We know that vampires are only $0.1\%$ of the population. 
 
 $\bullet$ $P(\text{vampire})=0.001$

}

\pause

\begin{center}
 Someone receive a positive test: Is she a vampire?
\end{center}
\pause
\begin{equation*}
 P(\text{vampire}|\text{positive}) = \frac{P(\text{positive}|\text{vampire})P(\text{vampire})}{P(\text{positive})}
\end{equation*}

\pause

\begin{center}
 She has \textbf{only 8.7\% chance} to actually be a vampire!
\end{center}




\end{frame}

\subsection{The inferential jump}

\begin{frame}
\begin{textblock}{128}(0,8)
\begin{center}
 \Large The inferential jump
 
 \normalsize (or inverse probability)
\end{center}
\end{textblock}
\vspace{0.9cm}

 \begin{framed}
 \centering
  All elements of previous examples were frequencies of observations
  
 \textbf{Bayesian inference quantify uncertainty about hidden variables}

 Therefore, we need \textbf{belief distributions} for those hidden variables!
 \end{framed} 

 \pause 

 \begin{equation}
  \underbrace{P(\text{Belief}|\text{Data})}_{\text{\scriptsize Posterior}} = \frac{\overbrace{P(\text{Data}|\text{Belief})}^{\text{\scriptsize Likelihood}} \overbrace{P(\text{Belief})}^{\text{\scriptsize Prior}} }{\underbrace{P(\text{Data})}_{\hfrac{\text{\scriptsize Evidence or}}{\text{\scriptsize Average likelihood}}}}
 \end{equation}

 \pause \vspace{0.3cm}
 
 
 \centering{In order to update our beliefs after seeing new data (posterior),
  we need to consider every possible path that could led as the data (likelihood).}
 
\end{frame}

\section{The inferential jump}

\subsection{The garden of forking paths}


\begin{frame}
\begin{textblock}{128}(0,8)
\begin{center}
 \large The garden of forking paths
\end{center}
\end{textblock}
\vspace{0.5cm}

\begin{textblock}{128}(0,21)
\centering
 Data: \includegraphics[page=1,width=0.065\textwidth]{imagenes/forkingPath} \hspace{0.3cm}
 Beliefs: \includegraphics[page=2,width=0.075\textwidth]{imagenes/forkingPath}, \includegraphics[page=3,width=0.075\textwidth]{imagenes/forkingPath}, \includegraphics[page=4,width=0.075\textwidth]{imagenes/forkingPath}, \includegraphics[page=5,width=0.075\textwidth]{imagenes/forkingPath}, \includegraphics[page=6,width=0.075\textwidth]{imagenes/forkingPath} 
 
\end{textblock}

\only<1>{
\begin{textblock}{128}(0,28)
\centering
Model:\vspace{0.1cm}

\includegraphics[page=12,width=0.6\textwidth]{imagenes/forkingPath}

\end{textblock}
}


\only<2>{
\begin{textblock}{128}(0,28)
\centering
Model:\vspace{0.1cm}

\includegraphics[page=7,width=0.6\textwidth]{imagenes/forkingPath}

\scriptsize Belief: \includegraphics[page=2,width=0.05\textwidth]{imagenes/forkingPath}
\end{textblock}
}


\only<3>{
\begin{textblock}{128}(0,28)
\centering
Model:\vspace{0.1cm}

\includegraphics[page=8,width=0.6\textwidth]{imagenes/forkingPath}

\scriptsize Belief: \includegraphics[page=3,width=0.05\textwidth]{imagenes/forkingPath}
\end{textblock}
}

\only<4>{
\begin{textblock}{128}(0,28)
\centering
Model:\vspace{0.1cm}

\includegraphics[page=9,width=0.6\textwidth]{imagenes/forkingPath}

\scriptsize Belief: \includegraphics[page=4,width=0.05\textwidth]{imagenes/forkingPath}
\end{textblock}
}

\only<5>{
\begin{textblock}{128}(0,28)
\centering
Model:\vspace{0.1cm}

\includegraphics[page=10,width=0.6\textwidth]{imagenes/forkingPath}

\scriptsize Belief: \includegraphics[page=5,width=0.05\textwidth]{imagenes/forkingPath}
\end{textblock}
}

\only<6>{
\begin{textblock}{128}(0,28)
\centering
Model:\vspace{0.1cm}

\includegraphics[page=11,width=0.6\textwidth]{imagenes/forkingPath}

\scriptsize Belief: \includegraphics[page=6,width=0.05\textwidth]{imagenes/forkingPath}
\end{textblock}
}

\only<2>{
\begin{textblock}{128}(0,70)
 \centering \scriptsize
 
 \begin{table}[H]
\begin{tabular}{cc}
 \ \ \ Belief \ \ \ & Ways to produce \includegraphics[page=1,width=0.045\textwidth]{imagenes/forkingPath}   \\ \hline
  \includegraphics[page=2,width=0.05\textwidth]{imagenes/forkingPath} &  $0 \times 4 \times 0 = 0$ \\
 & \\
 & \\
 & \\
 & \\
\end{tabular}
\end{table}
\end{textblock}

}

\only<3>{
\begin{textblock}{128}(0,70)
 \centering \scriptsize
 
 \begin{table}[H]
\begin{tabular}{cc}
 \ \ \ Belief \ \ \ & Ways to produce \includegraphics[page=1,width=0.045\textwidth]{imagenes/forkingPath}   \\ \hline
  \includegraphics[page=2,width=0.05\textwidth]{imagenes/forkingPath} &  $0 \times 4 \times 0 = 0$ \\
 \includegraphics[page=3,width=0.05\textwidth]{imagenes/forkingPath} &  $1 \times 3 \times 1 = 3$ \\
 & \\
 & \\
 & \\
\end{tabular}
\end{table}
\end{textblock}

}

\only<4>{
\begin{textblock}{128}(0,70)
 \centering \scriptsize
 
 \begin{table}[H]
\begin{tabular}{cc}
 \ \ \ Belief \ \ \ & Ways to produce \includegraphics[page=1,width=0.045\textwidth]{imagenes/forkingPath}   \\ \hline
  \includegraphics[page=2,width=0.05\textwidth]{imagenes/forkingPath} &  $0 \times 4 \times 0 = 0$ \\
 \includegraphics[page=3,width=0.05\textwidth]{imagenes/forkingPath} &  $1 \times 3 \times 1 = 3$ \\
 \includegraphics[page=4,width=0.05\textwidth]{imagenes/forkingPath} &  $2 \times 2 \times 2 = 8$ \\
 & \\
 & \\
\end{tabular}
\end{table}
\end{textblock}

}

\only<5>{
\begin{textblock}{128}(0,70)
 \centering \scriptsize
 
 \begin{table}[H]
\begin{tabular}{cc}
 \ \ \ Belief \ \ \ & Ways to produce \includegraphics[page=1,width=0.045\textwidth]{imagenes/forkingPath}   \\ \hline
  \includegraphics[page=2,width=0.05\textwidth]{imagenes/forkingPath} &  $0 \times 4 \times 0 = 0$ \\
 \includegraphics[page=3,width=0.05\textwidth]{imagenes/forkingPath} &  $1 \times 3 \times 1 = 3$ \\
 \includegraphics[page=4,width=0.05\textwidth]{imagenes/forkingPath} &  $2 \times 2 \times 2 = 8$ \\
 \includegraphics[page=5,width=0.05\textwidth]{imagenes/forkingPath} &  $3 \times 1 \times 3 = 9$ \\
 & \\
\end{tabular}
\end{table}
\end{textblock}

}

\only<6>{
\begin{textblock}{128}(0,70)
 \centering \scriptsize
 
 \begin{table}[H]
\begin{tabular}{cc}
 \ \ \ Belief \ \ \ & Ways to produce \includegraphics[page=1,width=0.045\textwidth]{imagenes/forkingPath}   \\ \hline
  \includegraphics[page=2,width=0.05\textwidth]{imagenes/forkingPath} &  $0 \times 4 \times 0 = 0$ \\
 \includegraphics[page=3,width=0.05\textwidth]{imagenes/forkingPath} &  $1 \times 3 \times 1 = 3$ \\
 \includegraphics[page=4,width=0.05\textwidth]{imagenes/forkingPath} &  $2 \times 2 \times 2 = 8$ \\
 \includegraphics[page=5,width=0.05\textwidth]{imagenes/forkingPath} &  $3 \times 1 \times 3 = 9$ \\
 \includegraphics[page=6,width=0.05\textwidth]{imagenes/forkingPath} &  $4 \times 0 \times 4 = 0$ \\
\end{tabular}
\end{table}
\end{textblock}

}

\end{frame}

\begin{frame}
 
 \footnotesize
 
 $\bullet$ \textbf{Prior} belief (distribution):
 \begin{equation*}
P(B) = \frac{1}{\# \text{Beliefs} }  \ \ \ \ \  \forall B \in \text{Beliefs}  
 \end{equation*}

\pause \vspace{0.3cm}
 
 $\bullet$ \textbf{Likelihood} (distribution):
 \begin{equation*}
P(D|B) = \frac{\text{Ways to produce $D$ given $B$}}{\text{Total ways given $B$}} \ \ \ \ \  \forall B \in \text{Beliefs}  
 \end{equation*}

\pause \vspace{0.3cm}

$\bullet$ \textbf{Evidence} or Average likelihood (scalar):
\begin{equation*}
 P(D) = \sum_{B \in \text{Beliefs}} \underbrace{P(D|B)}_{\text{likelihood}} \underbrace{P(B)}_{\text{prior}}
\end{equation*}

\pause \vspace{0.3cm}

$\bullet$ \textbf{Posterior} belief (distribution):

\begin{equation*}
 P(B|D) = \frac{P(D|B)P(D)}{P(D)}\ \ \ \ \  \forall B \in \text{Beliefs}  
\end{equation*}

 
 
\end{frame}

\subsection{The garden of (continuous) forking paths }

\begin{frame}
 \begin{textblock}{128}(0,8)
\begin{center}
 \large The garden of (continuous) forking paths 
\end{center}
\end{textblock}
\vspace{1.25cm}


\centering{How to estimate skill of players?}


  \begin{figure}[H]     
     \centering \normalsize
     \begin{subfigure}[b]{0.4\textwidth}
       \includegraphics[page=1,width=\textwidth]{../../Imagenes/elo.jpeg} 
       \caption*{Arpad Elo}
     \end{subfigure}
\end{figure}

\end{frame}

\begin{frame}
 \begin{textblock}{128}(0,8)
\begin{center}
 \normalsize Elo model
\end{center}
\end{textblock}

\centering{
\includegraphics[width=0.9\textwidth]{imagenes/elo-en}   
}

\end{frame}

\begin{frame}
 \begin{textblock}{128}(0,8)
\begin{center}
 \normalsize TrueSkill
\end{center}
\end{textblock}
\vspace{0.5cm}

\centering{
\includegraphics[width=1\textwidth]{imagenes/trueskill_versionGeneral-en}   
}
 
 
\begin{equation*}
 P(s|o,A) = \frac{P(o|s,A)P(s)}{P(o|A)}
\end{equation*}

 
\end{frame}

\begin{frame}
 \begin{textblock}{128}(0,8)
\begin{center}
 \normalsize Posterior
\end{center}
\end{textblock}
\vspace{0.5cm}

\Wider[4em]{ 
\begin{equation*}
 P(s_i \mid o, A) \propto 
 \begin{cases}
  N(s_i \, | \, \mu_i, \sigma_i^2) \, \Phi(\delta(s_i) \, | \, 0, \vartheta - \sigma_i^2) & \text{Win} \\
  N(s_i \, | \, \mu_i, \sigma_i^2 ) \, (1 - \Phi(\delta(s_i) \, | \, 0, \vartheta - \sigma_i^2)) & \text{Loose}
 \end{cases}
\end{equation*}

\includegraphics[width=0.49\textwidth]{imagenes/posterior_win}   
\includegraphics[width=0.49\textwidth]{imagenes/posterior_loose}   
}
\end{frame}


\section{Bayesian model inference}

\begin{frame}
\begin{textblock}{128}(0,8)
\begin{center}
 \Large Bayesian model inference
\end{center}
\end{textblock}
\vspace{0.5cm}
\Wider[3em]{ 
$\bullet$ Which are our beliefs about different hidden models $M$?
\pause
\begin{equation*}
 P(M_k|D) = \frac{P(D|M_k)P(M_k)}{\sum_{i=1}^n P(D|M_i) P(M_i)}
\end{equation*}

\pause \vspace{0.2cm}

$\bullet$To compare models we can compute their ratio probability,
\begin{equation*}
 \text{Bayes factor}(q,r) = \frac{P(M_q|D)}{P(M_r|D)} = \frac{P(D|M_q)P(M_q)}{P(D|M_r)P(M_r)} \pause = \frac{P(D|M_q)}{\underbrace{P(D|M_r)}_{\text{\scriptsize Evidence!}}}   
\end{equation*}

\pause \vspace{0.2cm}

}

\begin{mdframed}[backgroundcolor=black!20]
\centering
$P(M_q|D) > P(M_r|D) \Longleftrightarrow P(D|M_q) > P(D|M_r)$

\end{mdframed}


\begin{center}
 \Large All you need is evidence
\end{center} 

 
\end{frame}

\subsection{Evidence}

\begin{frame}
\begin{textblock}{128}(0,8)
\begin{center}
 \large Evidence
\end{center}
\end{textblock}
\vspace{1cm}
 
 \centering{
 \includegraphics[width=1\textwidth]{imagenes/evidence}
 }
 
\end{frame}


\begin{frame}
\begin{textblock}{128}(0,8)
\begin{center}
 \large Example
\end{center}
\end{textblock}
\vspace{0.75cm}


\Wider[1cm]{
\begin{center}
 Which model has highest probability given the observed data $D$
 \end{center}
}


 \begin{table}[H]
\begin{tabular}{|c|c|c|c|c|c|c|c|c|c|c|}
\hline
 Ana & 0.6 & 0.4 & 0.4 & 0.4 & 0.6 & 0.6 & 0.6 & 0.6 & 0.6 & 0.6 \\ 
\hline \hline
 Observed& \includegraphics[width=0.02\textwidth]{imagenes/sol}  & \includegraphics[width=0.02\textwidth]{imagenes/lluvia} & \includegraphics[width=0.02\textwidth]{imagenes/lluvia} & \includegraphics[width=0.02\textwidth]{imagenes/lluvia} & \includegraphics[width=0.02\textwidth]{imagenes/sol} & \includegraphics[width=0.02\textwidth]{imagenes/sol} & \includegraphics[width=0.02\textwidth]{imagenes/sol} & \includegraphics[width=0.02\textwidth]{imagenes/sol} & \includegraphics[width=0.02\textwidth]{imagenes/sol} & \includegraphics[width=0.02\textwidth]{imagenes/sol} \\
\hline \hline
 Berta & 0 & 0 & 1 & 1 & 1 & 0 & 0 & 0 & 0 & 0   \\
 \hline
\end{tabular}
\end{table}

\pause

\Wider[1cm]{
\begin{itemize}
 \item[$\bullet$] Berta seems to do a better job according to all data science metrics
 
 \item[$\bullet$] However Ana's model assigns highest probability to the observed data!! 
\end{itemize}
}

\pause \vspace{0.2cm}

  \begin{table}[H]
\begin{tabular}{c|c}
 & $P(D|M) = \prod P(D_i|M)$ \\ \hline
 Ana &  $0.4^{10}= 10.1E^{-5}$  \\
Berta & $1.0^8  \times 0^2 = 0 = 1E^{-\infty}$\\
 \end{tabular}
\end{table}
\vspace{-0.1cm}
\begin{framed}
\centering 

Ana has an infinitely more likely model
\end{framed}


\end{frame}


% 
% \subsection{Metricas}
% 
% \begin{frame}
% \begin{table}[H]
% \centering
% \begin{tabular}{rlll}
%   \hline
%  & Ana & Berta & Gana \\ 
%   \hline
% Mean squared error & 3.6 & 2 & Berta \\ 
%   Mean absolute error & 6 & 2 & Berta \\ 
%   Hit rate & 0.4 & 0.8 & Berta \\ 
%   Accuracy & 0 & 0.8 & Berta \\ 
%   Precsion & 0 & 0.67 & Berta \\ 
%   Recall & NaN & 0.67 & Berta \\ 
%   F1 score & NaN & 0.25 & Berta \\ 
%   ROC area & 0 & 0.76 & Berta \\ 
%   Evidence & $10.4E^{-5}$ & 0 & Ana \\ 
%    \hline
% \end{tabular}
% \end{table}
% 
% 
% \end{frame}
% 




\begin{frame}
\begin{textblock}{128}(0,8)
\begin{center}
 \large Linear regression example
\end{center}
\end{textblock}
\vspace{0.75cm}

\centering{
 \includegraphics[width=0.8\textwidth]{imagenes/minka.png}
 }


\end{frame}


\begin{frame}
\begin{textblock}{128}(0,8)
\begin{center}
 \normalsize Evidence vs maximum likelihood
\end{center}
\end{textblock}
\vspace{1cm}
\Wider[4em]{

\centering{
 \includegraphics[width=0.49\textwidth]{imagenes/minka_evidence.png}
 \includegraphics[width=0.49\textwidth]{imagenes/minka_likelihood.png}
 }
}

\begin{mdframed}[backgroundcolor=black!15]
\centering \Large
We do not need any test set
\end{mdframed}





\end{frame}


\section{Bayesian no-doubt case}

\begin{frame}
 \begin{textblock}{128}(0,8)
\begin{center}
 \Large Bayesian no-doubt case
\end{center}
\end{textblock}
\vspace{0.5cm}


 \centering{Fixed beliefs, even with infinite new data}
 \begin{mdframed}[backgroundcolor=black!15]
 \centering
 $\#$Beliefs = 1 $\Longrightarrow$ $\underbrace{P(B)}_{\text{\scriptsize Prior}} = \underbrace{P(B|D)}_{\text{\scriptsize Posterior}} \ \ \ \ \hfrac{\forall D \in \text{Data}}{\forall B \in \text{Beliefs}}$
\end{mdframed}

 \vspace{0.9cm}
\pause 


 \centering{Likelihood is just the Evidence}
\begin{mdframed}[backgroundcolor=black!15]
\centering
 $\#$Beliefs = 1  $\Longleftrightarrow$ Likelihood = Evidence  
\end{mdframed}
 

 

\end{frame}



\subsection{The hacked-belief approach}

\begin{frame}
 \begin{textblock}{128}(0,8)
\begin{center}
 \large The hacked-belief approach
\end{center}
\end{textblock}
\vspace{0.5cm}


\begin{center}
 
 \begin{mdframed}
 \centering
  Who has no doubt? Who has only one belief?
 \end{mdframed}
 
 \vspace{0.3cm}
 \pause
 
 \begin{itemize}
  \item[$\bullet$] God (if exists) \pause
  \item[$\bullet$] Mathematicians (and others non empricial sciences) \pause
  \item[$\bullet$] Maybe some extremists \pause
  \item[$\bullet$] \huge All non-bayesian machine learning
 \end{itemize}

 
 
 
\end{center}

 
\end{frame}


\begin{frame}
 \begin{textblock}{128}(0,8)
\begin{center}
 \large The great dishonesty
\end{center}
\end{textblock}

\begin{mdframed}
 \centering{ $\overset{\text{The best belief after seeing the data}}{\text{\st{maximum likelihood estimator} }}$ = $\underset{B }{\text{argmax}} P(D|B,M) = \widehat{B}_{D,M}$  }
\end{mdframed}

\vspace{0.6cm}

 \begin{equation*}
\begin{split}
P(D|\widehat{B}_{D,M},M) & = \overbrace{P(D|\underset{B}{\text{argmax}} P(D|B,M), M)}^{\hfrac{\text{\scriptsize Hacked \textbf{likelihood} }}{\text{\scriptsize  given training data}}} = \overbrace{P(D|D, M)}^{\hfrac{\text{\scriptsize Hacked \textbf{evidence} }}{\text{\scriptsize given training data}}}%\\
%&= \int_B P(D|B,M) \mathbb{I}(B=\widehat{B}_{D,M} ) dB
\end{split}
 \end{equation*}


 \pause
 \vspace{0.6cm}
 
\centering{\large With great dishonesty, comes great regularization {\tiny Zivik talk} }
 
 

\end{frame}



% 
% 
% 
% \begin{frame}
%  
%  \begin{center}
%   \Large Annex
%  \end{center}
% 
%  
% \end{frame}
% 
% 
% \begin{frame}
% \begin{textblock}{128}(0,8)
% \begin{center}
% Annex 1 \\ 
%  \scriptsize Proof: Deviance -- Bayes factor
% \end{center}
% \end{textblock}
% \vspace{0.5cm}
% 
% 
% \tiny
% 
% % 
% Entropy of a probability distribution
% \begin{equation}
%  H(P(C)) = -E(log(P(C))) = - \sum_{i=1}^n P(C_i) log(P(C_i))
% \end{equation}
% 
% \pause
% \vspace{0.3cm}
% 
% $\bullet$ The log Bayes factor is essentially the relative distance to the truth
% 
% \begin{mdframed}[backgroundcolor=black!15]
% \begin{equation}\footnotesize
%   D_{KL}(p,q) - D_{KL}(p,r) \overset{\text{Annex 1}}{=} \frac{1}{|D|} \text{log}\left(\text{Bayes factor}(r,q)\right) 
% \end{equation}
% \end{mdframed}
% 
% 
% 
% \begin{equation}
%  D_{KL}(p,q) = \sum_{i=1}^{|C|} p_i  (\text{log}(p_i) - \text{log}(q_i))
% \end{equation}
% With truth $P(C_i)=p_i$ and model $P(C_i|M_q)=q_i$, for each event $C_i \in C$%, and
% 
% 
% 
% 
% \begin{equation}
% \text{Deviance}(q,r) = D_{KL}(p,q) - D_{KL}(p,r) = \underbrace{\sum_{i=1}^{|C|} p_i \text{log}(r_i)}_{E[\text{log}(r)]} -\underbrace{\sum_{i=1}^{|C|} p_i \text{log}(q_i)}_{E[\text{log}(q)]}   
%  \end{equation}
% 
%  
%  \vspace{0.3cm}
%  
%  Where
%  \begin{equation*}
% \begin{align*}
%  E[\text{log}(q)] & = \sum_{i=1}^{|C|} P(C_i) \text{log}(P(C_i|M_q))  \approx  \frac{1}{|D|} \sum_{i=1}^{|D|} \text{log}(P(D_i|M_q)) \\
%  & = \frac{1}{|D|} \text{log}\left(\prod_{i=1}^{|D|} P(D_i|M_q)\right)= \frac{1}{|D|} \text{log}\underbrace{(P(D|M_q))}_{\text{Evidence!}}
% \end{align*}
% \end{equation*}
% 
% \vspace{0.3cm}
% 
% Then
% \begin{equation*}
% \begin{align*}
% D_{KL}(p,q) - D_{KL}(p,r) & \approx \frac{1}{|D|} \text{log}(P(D|M_q)) - \frac{1}{|D|} \text{log}(P(D|M_q)) \\
% & = \frac{1}{|D|} \text{log}\left(\frac{P(D|M_r)}{P(D|M_q)} \right) \\
% & = \frac{1}{|D|} \text{log}\left(\text{Bayes factor}\right) 
% \end{align*}
%  \end{equation*}
%  
% \end{frame}

\end{document}



